\documentclass[12pt, a4paper]{article}
\usepackage[utf8]{inputenc}
\usepackage[T1]{fontenc}
\usepackage{courier}
\usepackage[french]{babel}
\usepackage{graphicx} % Required for inserting images
\usepackage{amsmath}
\usepackage{amssymb}
\usepackage{amsthm}
\usepackage{bbold}
\usepackage{stmaryrd}
\usepackage[hidelinks]{hyperref}
\usepackage[margin=2.5cm]{geometry}
\usepackage{listings}
\usepackage{pgfplots}
\usepackage{caption}
\usepackage{subcaption}
\usepackage{tikz}
\usetikzlibrary{positioning,intersections}
\usepackage{tikzit}
\input{sample.tikzstyles}
\input{quantum-circuits.tikzstyles}
\input{quantum-circuits.tikzdefs}

\pgfplotsset{compat=1.18}

\lstdefinelanguage{ExampleIsar}%
{morekeywords={
_first,theory,end,types,datatype,consts,defs,primrec,
syntax,translations,apply, lemma,theorem,corollary,
done,sorry,goal,proof,fixed variables,_last,simproc_setup},
sensitive=true,
morecomment=[l]{-- },
literate=
{\\<not>}{{$¬$}}1 {\\<times>}{{$×$}}1
{\\<Rightarrow>}{{$\Rightarrow$}}1%
{\\<equiv>}{{$\equiv$}}1 {~=}{{$\not=$}}1
{\\<rightleftharpoons>}
{{$\rightleftharpoons$}}1
{\\<exists>}{{$\exists$}}1
{\\<parallel>}{{$\parallel$}}1
{\\<Longrightarrow>}
{{$\Longrightarrow$}}1
{\\<lbrakk>}{{$\lsemantics$}}1
{\\<rbrakk>}{{$\rsemantics$}}1
{\\<longrightarrow>}
{{$\longrightarrow$}}1
{\\<and>}{{$\land$}}1
{\\<or>}{{$\lor$}}1
{\\<bigwedge>}{{$\bigwedge$}}1
{ü}{"u}1 {ä}{"a}1 {ö}{"o}1
{Ü}{"U}1 {Ä}{"A}1 {Ö}{"O}1 {ß}{{ss}}1
}[keywords,comments,strings]%

\definecolor{keyword1}{rgb}{0.2,0.3,0.6}
\definecolor{keyword2}{rgb}{0.3,0.6,0.2}
\definecolor{keyword3}{rgb}{1,0,0}
\definecolor{keyword4}{rgb}{0.8,0.5,0.9}
\definecolor{keyword5}{rgb}{0,0.8,1}

\lstdefinelanguage{Isar}{
  basicstyle=\ttfamily\small,
  sensitive = true,
  keywords = {session,theory,lemma,lemmas,bundle,datatype,function,fun,inductive,simproc_setup,have,including,then,also,moreover,finally,ultimately,using,by,proof,qed},
  keywords = [2]{sessions,directories,theories,imports,begin,end,where,passive,weak_congproc},
  keywords = [3]{t,t1,t2,u,u1,u2,v,v1,v2,x,x1,x2,y,y1,y2},
  keywords = [4]{apply},
  keywords = [5]{add,only,cong},
  keywords = [6]{show},
  keywordstyle=\color{keyword1}\bfseries,
  keywordstyle=[2]\color{keyword2}\bfseries,
  keywordstyle=[3]\color{blue},
  keywordstyle=[4]\color{keyword3}\bfseries,
  keywordstyle=[5]\color{keyword4},
  keywordstyle=[6]\color{keyword5}\bfseries,
  otherkeywords={% Operators
    >, <, ==
  },
  showstringspaces=false,
  breaklines=true,
  morecomment=[s]{(*}{*)},
  commentstyle=\color{red}\ttfamily,
  literate=
    {\\<open>}{\guilsinglleft}1
    {\\<close>}{\guilsinglright}1
    {\\<lbrakk>}{$\llbracket$}1
    {\\<rbrakk>}{$\rrbracket$}1
}

\newtheorem{definition}{Définition}[subsection]
\newtheorem{theorem}[definition]{Théorème}
\newtheorem{corollary}[definition]{Corollaire}
\newtheorem{lemma}[definition]{Lemme}
\newtheorem{property}[definition]{Propriété}
\newtheorem{remark}[definition]{Remarque}
\newtheorem{example}[definition]{Exemple}

\title{Certificats pour les preuves d'équivalence de circuits avec Isabelle/HOL}
\date{29 mai 2026}
\author{Alexis Beaucourt}

\newcommand{\tinyinput}[1]{\scalebox{0.44}{\input{#1}}}
\newcommand{\layer}[3]{L^{\smash{#1}}_{\smash{#2}}(#3)}

\def\prob#1{ \mathbb P \left( #1 \right) }
\def\esp#1{ \mathbb E \left[ #1 \right] }
\def\var#1{ \text{Var} \left( #1 \right) }

\def\N{ \mathbb N }
\def\R{ \mathbb R }

\def\setcomp#1#2{ \left\{ #1 \, \middle| \, #2 \right\} }
\def\e{\varepsilon}
\def\leqv{\Longleftrightarrow}
\def\id{ \mathrm{id} }
\def\nf{ \mathrm{NF} }
\def\pro#1{ \mathbf{PRO}_{ #1 } }
\def\prop#1{ \mathbf{PROP}_{ #1 } }
\def\dom{ \mathrm{dom} }
\def\cod{ \mathrm{cod} }
\def\clr#1{ \overset{=}{ #1 } }
\def\clt#1{ \overset{+}{ #1 } }
\def\clrt#1{ \overset{*}{ #1 } }

% abstract potentiel
% \tableofcontents
% intro (etat de l'art & motivation & description contenu)
% conclusion
% remerciments (loria, veridis, encadrants, g5k)

\begin{document}

\maketitle


\section{Introduction}

\section{Définition des PRO/PROP}

Chaque circuit $u$ possède un domaine $\dom(u) \in \N$ et un codomaine $\cod(u) \in \N$

On définit les circuits de $\pro{\Sigma, \mathcal E}$ comme étant les circuits formés à partir des générateurs suivants:

\begin{tikzpicture}
	\begin{pgfonlayer}{nodelayer}
		\node [style=none] (0) at (-3, 0.5) {};
		\node [style=none] (1) at (-1, 0.5) {};
		\node [style=none] (2) at (-3, -0.5) {};
		\node [style=none] (3) at (-1, -0.5) {};
		\node [style=none] (4) at (-2, 0) {$\vdots$};
		\node [style=none] (5) at (-2, -1.25) {$\id_n$};
		\node [style=none] (6) at (-2, -1.75) {$n \in \N$};
		\node [style=none] (7) at (0.75, 0.5) {};
		\node [style=none] (8) at (2.25, 0.5) {};
		\node [style=none] (9) at (2.75, 0.5) {};
		\node [style=none] (10) at (4.25, 0.5) {};
		\node [style=none] (11) at (0.75, -0.5) {};
		\node [style=none] (12) at (2.25, -0.5) {};
		\node [style=none] (13) at (2.75, -0.5) {};
		\node [style=none] (14) at (4.25, -0.5) {};
		\node [style=medium box] (15) at (2.5, 0) {$u$};
		\node [style=none] (16) at (2.5, -1.75) {$u \in \Sigma$};
		\node [style=none] (17) at (1.5, 0) {$\vdots$};
		\node [style=none] (18) at (3.25, 0) {$\vdots$};
		\node [style=none] (19) at (1.5, -1) {$\dom(u)$};
		\node [style=none] (20) at (3.5, -1) {$\cod(u)$};
	\end{pgfonlayer}
	\begin{pgfonlayer}{edgelayer}
		\draw (0.center) to (1.center);
		\draw (2.center) to (3.center);
		\draw (7.center) to (8.center);
		\draw (9.center) to (10.center);
		\draw (11.center) to (12.center);
		\draw (13.center) to (14.center);
		\draw (19.center) to (19.center);
	\end{pgfonlayer}
\end{tikzpicture}


Ainsi que les opérations suivantes:

\input{pic/ops.tikz}

Cependant, on quotiente le tout par les égalités de $\mathcal E$, en plus des égalités suivantes:


\begin{figure}[h]
    \fbox{\begin{minipage}{0.982\linewidth}\begin{center}
        \vspace{-1em}
        \hspace{-0.8em}
        \begin{subfigure}{0.48\textwidth}
            \begin{equation}\label{eq:seqidentity}
                \begin{array}{ccccc}
                    \id_n\fatsemi t&=&t&=&t\fatsemi\id_m\\[.4em]
                    \tinyinput{fig/seqid_left.tikz}&=&\tinyinput{fig/seqid_mid.tikz}&=&\tinyinput{fig/seqid_right.tikz}
                \end{array}
            \end{equation}
        \end{subfigure}

        \vspace{-0.5em}
        \begin{subfigure}{0.44\textwidth}
            \begin{equation}\label{eq:seqassociativity}
                \begin{array}{ccc}
                    (t\fatsemi u)\fatsemi v&=&t\fatsemi(u\fatsemi v)\\[.4em]
                    \tinyinput{fig/seqasso_left.tikz}&=&\tinyinput{fig/seqasso_right.tikz}
                \end{array}
            \end{equation}
        \end{subfigure}\hspace{2em}
        \begin{subfigure}{0.38\textwidth}
            \begin{equation}\label{eq:parassociativity}
                \begin{array}{ccc}
                    (t\otimes u)\otimes v&=&t\otimes(u\otimes v)\\[.4em]
                    \tinyinput{fig/parasso_left.tikz}&=&\tinyinput{fig/parasso_right.tikz}
                \end{array}
            \end{equation}
        \end{subfigure}

        \vspace{-0.5em}
        \begin{subfigure}{0.30\textwidth}
            \begin{equation}\label{eq:idinduct}
                \begin{array}{ccc}
                    \id_{n+m}&=&\id_{n}\otimes\id_{m}\\[.4em]
                    \tinyinput{fig/idpar_left.tikz}&=&\tinyinput{fig/idpar_right.tikz}
                \end{array}
            \end{equation}
        \end{subfigure}\hspace{2em}
        \begin{subfigure}{0.40\textwidth}
            \begin{equation}\label{eq:paridentity}
                \begin{array}{ccccc}
                    t\otimes \id_0&=&t&=&\id_0\otimes t\\[.4em]
                    \tinyinput{fig/parid_left.tikz}&=&\tinyinput{fig/parid_mid.tikz}&=&\tinyinput{fig/parid_right.tikz}
                \end{array}
            \end{equation}
        \end{subfigure}

        \vspace{-0.5em}
        \begin{subfigure}{0.57\textwidth}
            \begin{equation}\label{eq:interchange}
                \begin{array}{ccc}
                    (u_1\fatsemi u_2)\otimes(v_1\fatsemi v_2)&=&(u_1\otimes v_1)\fatsemi(u_2\otimes v_2)\\[.4em]
                    \tinyinput{fig/interchange_left.tikz}&=&\tinyinput{fig/interchange_right.tikz}
                \end{array}
            \end{equation}
        \end{subfigure}\hspace{1em}
        \begin{subfigure}{0.36\textwidth}
            \begin{equation}\label{eq:emptysymmetry}
                \begin{array}{ccccc}
                    \sigma_{n,0}&=&\id_n&=&\sigma_{0,n}\\[.4em]
                    \tinyinput{fig/s0_left.tikz}&=&\tinyinput{fig/s0_mid.tikz}&=&\tinyinput{fig/s0_right.tikz}
                \end{array}
            \end{equation}
        \end{subfigure}

        \vspace{-0.5em}
        \begin{subfigure}{0.34\textwidth}
            \begin{equation}\label{eq:involution}
                \begin{array}{ccc}
                    \sigma_{n,m}\fatsemi\sigma_{m,n}&=&\id_{n+m}\\[.4em]
                    \tinyinput{fig/inv_left.tikz}&=&\tinyinput{fig/inv_right.tikz}
                \end{array}
            \end{equation}
        \end{subfigure}\hspace{2em}
        \begin{subfigure}{0.48\textwidth}
            \begin{equation}\label{eq:naturality}
                \begin{array}{ccc}
                    (t\otimes\id_k)\fatsemi\sigma_{m,k}&=&\sigma_{n,k}\fatsemi(\id_k\otimes t)\\[.4em]
                    \tinyinput{fig/pass_left.tikz}&=&\tinyinput{fig/pass_right.tikz}
                \end{array}
            \end{equation}
        \end{subfigure}

        \vspace{-0.5em}
        \begin{subfigure}{0.86\textwidth}
            \begin{equation}\label{eq:symmetryinduct}
                \begin{array}{ccc}
                    \sigma_{n+k,m+\ell}&=&(\id_n\otimes\sigma_{k,m}\otimes\id_\ell)\fatsemi(\sigma_{n,m}\otimes\sigma_{k,\ell})\fatsemi(\id_m\otimes\sigma_{n,\ell}\otimes\id_k)\\[.4em]
                    \tinyinput{fig/swap_decomp_left.tikz}&=&\tinyinput{fig/swap_decomp_right.tikz}
                \end{array}
            \end{equation}
        \end{subfigure}
        \vspace{0.2em}
    \end{center}\end{minipage}}
\end{figure}



On peut ensuite définir les $\prop{\Sigma, \mathcal E}$ comme étant des $\pro{\Sigma \cup \Sigma_0, \mathcal E \cup \mathcal E_0}$, où $\Sigma_0$ contient:

\begin{tikzpicture}
	\begin{pgfonlayer}{nodelayer}
		\node [style=none] (0) at (-1.5, 1.75) {};
		\node [style=none] (1) at (-1.5, 0.75) {};
		\node [style=none] (2) at (-0.5, 1.75) {};
		\node [style=none] (3) at (-0.5, 0.75) {};
		\node [style=none] (4) at (-1.5, 0.25) {};
		\node [style=none] (5) at (-0.5, 0.25) {};
		\node [style=none] (6) at (-1.5, -0.75) {};
		\node [style=none] (7) at (-0.5, -0.75) {};
		\node [style=none] (8) at (0.5, 0.75) {};
		\node [style=none] (9) at (0.5, 1.75) {};
		\node [style=none] (10) at (1.5, 1.75) {};
		\node [style=none] (11) at (1.5, 0.75) {};
		\node [style=none] (12) at (0.5, 0.25) {};
		\node [style=none] (13) at (1.5, 0.25) {};
		\node [style=none] (14) at (0.5, -0.75) {};
		\node [style=none] (15) at (1.5, -0.75) {};
		\node [style=none] (16) at (-0.5, 1.75) {};
		\node [style=none] (24) at (-1, 1.25) {$\vdots$};
		\node [style=none] (28) at (-1, -0.25) {$\vdots$};
		\node [style=none] (29) at (1, 1.25) {$\vdots$};
		\node [style=none] (30) at (1, -0.25) {$\vdots$};
		\node [style=none] (31) at (0, -1.5) {$\sigma_{n,m}$};
		\node [style=none] (32) at (-1.75, 1.25) {$n$};
		\node [style=none] (33) at (1.75, -0.25) {$n$};
		\node [style=none] (34) at (-1.75, -0.25) {$m$};
		\node [style=none] (35) at (1.75, 1.25) {$m$};
		\node [style=none] (36) at (0, -2) {$n,m \in \N$};
	\end{pgfonlayer}
	\begin{pgfonlayer}{edgelayer}
		\draw (0.center) to (2.center);
		\draw (8.center) to (8.center);
		\draw (8.center) to (11.center);
		\draw (2.center) to (12.center);
		\draw (12.center) to (13.center);
		\draw (14.center) to (15.center);
		\draw (1.center) to (3.center);
		\draw (3.center) to (14.center);
		\draw (5.center) to (9.center);
		\draw (7.center) to (8.center);
		\draw (4.center) to (5.center);
		\draw (6.center) to (7.center);
		\draw (9.center) to (10.center);
	\end{pgfonlayer}
\end{tikzpicture}


Et les $\mathcal E_0$ sont les égalités:

[insérer image]

Pour le restant de ce rapport, on notera $\simeq$ la relation d'équivalence entre deux termes.


Afin de décider l'équivalence $t_1 \simeq t_2$, il convient de définir une fonction de normalisation sur les circuits. Cette fonction est de plus exprimée sous forme de système de réécriture, dont on prouve qu'il est terminant et confluent.

Pour simplifier les notations, on pose $\layer{k}{\ell}{g} = \textup{id}_k \otimes g \otimes \textup{id}_\ell$

Ainsi que le défaut $\textup{def}(u) = \textup{dom}(u) - \textup{cod}(u)$

Pour les $\pro{\Sigma}$, on a le système suivant:

\begin{figure}[h]
    \fbox{\begin{minipage}{0.982\linewidth}\begin{center}
        \vspace{-1em}
        \hspace{-0.8em}
        \begin{subfigure}{0.37\textwidth}
            \begin{equation}\label{ru:seqasso}\tag{R1}
                (t\fatsemi u)\fatsemi v \longrightarrow t\fatsemi (u\fatsemi v)
            \end{equation}
        \end{subfigure}\hspace{2em}
        \begin{subfigure}{0.42\textwidth}
            \begin{equation}\label{ru:parasso}\tag{R2}
                (t\otimes u)\otimes v \longrightarrow t\otimes (u\otimes v)
            \end{equation}
        \end{subfigure}

        \vspace{-0.3em}
        \hspace{-0.8em}
        \begin{subfigure}{0.25\textwidth}
            \begin{equation}\label{ru:seqidleft}\tag{R3}
                \textup{id}_n\fatsemi t \longrightarrow t
            \end{equation}
        \end{subfigure}\hspace{2em}
        \begin{subfigure}{0.25\textwidth}
            \begin{equation}\label{ru:seqidright}\tag{R4}
                t\fatsemi\textup{id}_n \longrightarrow t
            \end{equation}
        \end{subfigure}

        \vspace{-0.3em}
        \hspace{-0.8em}
        \begin{subfigure}{0.36\textwidth}
            \begin{equation}\label{ru:idid}\tag{R5}
                \textup{id}_n\otimes\textup{id}_m \longrightarrow \textup{id}_{n+m}
            \end{equation}
        \end{subfigure}\hspace{2em}
        \begin{subfigure}{0.46\textwidth}
            \begin{equation}\label{ru:ididctx}\tag{R6}
                \textup{id}_n\otimes(\textup{id}_m\otimes t) \longrightarrow \textup{id}_{n+m}\otimes t
            \end{equation}
        \end{subfigure}

        \vspace{-0.3em}
        \hspace{-0.8em}
        \begin{subfigure}{0.82\textwidth}
            \begin{equation}\label{ru:par}\tag{R7}
                u\otimes v \longrightarrow (u\otimes\textup{id}_{\textup{dom}(v)})\fatsemi(\textup{id}_{\textup{cod}(u)}\otimes v)
                \;\;\text{if}\;\; u\ne\textup{id}_n\wedge v\ne\textup{id}_m
            \end{equation}
        \end{subfigure}

        \vspace{-0.3em}
        \hspace{-0.8em}
        \begin{subfigure}{0.54\textwidth}
            \begin{equation}\label{ru:parseqidleft}\tag{R8}
                \textup{id}_n\otimes(u\fatsemi v) \longrightarrow (\textup{id}_n\otimes u)\fatsemi(\textup{id}_n\otimes v)
            \end{equation}
        \end{subfigure}

        \vspace{-0.3em}
        \hspace{-0.8em}
        \begin{subfigure}{0.54\textwidth}
            \begin{equation}\label{ru:parseqidright}\tag{R9}
                (u\fatsemi v)\otimes\textup{id}_n \longrightarrow (u\otimes\textup{id}_n)\fatsemi(v\otimes\textup{id}_n)
            \end{equation}
        \end{subfigure}

        \vspace{-0.3em}
        \hspace{-0.8em}
        \begin{subfigure}{0.73\textwidth}
            \begin{equation}\label{ru:layercom}\tag{R10}
                \begin{array}{c}
                    \layer{k_1}{\ell_1}{g_1}\fatsemi\layer{k_2}{\ell_2}{g_2} \longrightarrow \layer{k_2}{\ell_2+\textup{def}(g_1)}{g_2}\fatsemi \layer{k_1-\textup{def}(g_2)}{\ell_1}{g_1} \\
                    \text{if}\;\; k_1\ge k_2+\textup{dom}(g_2)
                \end{array}
            \end{equation}
        \end{subfigure}

        \vspace{-0.3em}
        \hspace{-0.8em}
        \begin{subfigure}{0.85\textwidth}
            \begin{equation}\label{ru:layercomctx}\tag{R11}
                \begin{array}{c}
                    \layer{k_1}{\ell_1}{g_1}\fatsemi (\layer{k_2}{\ell_2}{g_2}\fatsemi t) \longrightarrow \layer{k_2}{\ell_2+\textup{def}(g_1)}{g_2}\fatsemi (\layer{k_1-\textup{def}(g_2)}{\ell_1}{g_1}\fatsemi t) \\
                    \text{if}\;\; k_1\ge k_2+\textup{dom}(g_2)
                \end{array}
            \end{equation}
        \end{subfigure}
        \vspace{0.2em}
    \end{center}\end{minipage}}
\end{figure}

On définit un sorte spéciale de swap $\tau_n = \sigma_{n,1}$

Pour les $\prop{\varnothing}$, on a le système composé des règles \verb|R1_NF| jusqu'à \verb|R9_NF| ainsi que les règles suivantes:

\begin{figure}[h]
    \fbox{\begin{minipage}{0.982\linewidth}\begin{center}
        \vspace{-1em}
        \hspace{-0.8em}
        \begin{subfigure}{0.26\textwidth}
            \begin{equation}\label{ru:swap:swap0left}\tag{R12}
                \sigma_{0,n} \longrightarrow \textup{id}_n
            \end{equation}
        \end{subfigure}\hspace{3em}
        \begin{subfigure}{0.26\textwidth}
            \begin{equation}\label{ru:swap:swap0right}\tag{R13}
                \sigma_{n,0} \longrightarrow \textup{id}_n
            \end{equation}
        \end{subfigure}

        \vspace{-0.5em}
        \hspace{-0.8em}
        \begin{subfigure}{0.69\textwidth}
            \begin{equation}\label{ru:swap:swapreduce}\tag{R14}
                \sigma_{n+1,m+2} \longrightarrow (\sigma_{n+1,m+1}\otimes\textup{id}_1)\fatsemi(\textup{id}_{m+1}\otimes\tau_{n+1})
            \end{equation}
        \end{subfigure}        

        \vspace{-0.5em}
        \hspace{-0.8em}
        \begin{subfigure}{0.66\textwidth}
            \begin{equation}\label{ru:swap:layernat}\tag{R15}
                \begin{array}{c}
                    \layer{k_1}{\ell_1}{\tau_{d_1}}\fatsemi\layer{k_2}{\ell_2}{\tau_{d_2}} \longrightarrow \layer{k_2}{\ell_2}{\tau_{d_2}}\fatsemi \layer{k_1+1}{\ell_1-1}{\tau_{d_1}} \\
                    \text{if}\;\; d_1\ge1,d_2\ge1 \;\;\text{and}\;\; k_2\le k_1<k_2+d_2-d_1
                \end{array}
            \end{equation}
        \end{subfigure}

        \vspace{-0.5em}
        \hspace{-0.8em}
        \begin{subfigure}{0.76\textwidth}
            \begin{equation}\label{ru:swap:layernatctx}\tag{R16}
                \begin{array}{c}
                    \layer{k_1}{\ell_1}{\tau_{d_1}}\fatsemi (\layer{k_2}{\ell_2}{\tau_{d_2}}\fatsemi t) \longrightarrow \layer{k_2}{\ell_2}{\tau_{d_2}}\fatsemi (\layer{k_1+1}{\ell_1-1}{\tau_{d_1}}\fatsemi t) \\
                    \text{if}\;\; d_1\ge1,d_2\ge1 \;\;\text{and}\;\; k_2\le k_1<k_2+d_2-d_1
                \end{array}
            \end{equation}
        \end{subfigure}

        \vspace{-0.5em}
        \hspace{-0.8em}
        \begin{subfigure}{0.84\textwidth}
            \begin{equation}\label{ru:swap:layercon}\tag{R17}
                \begin{array}{c}
                    \layer{k_1}{\ell_1}{\tau_{d_1}}\fatsemi\layer{k_2}{\ell_2}{\tau_{d_2}} \longrightarrow \layer{k_2}{\ell_2+1}{\tau_{d_2-1}}\fatsemi \layer{k_1+1}{\ell_1}{\tau_{d_1-1}} \\
                    \text{if}\;\; d_1\ge1,d_2\ge1 \;\;\text{and}\;\; k_2\le k_1\wedge k_2+d_2-d_1\le k_1<k_2+d_2
                \end{array}
            \end{equation}
        \end{subfigure}

        \vspace{-0.5em}
        \hspace{-0.8em}
        \begin{subfigure}{0.84\textwidth}
            \begin{equation}\label{ru:swap:layerconctx}\tag{R18}
                \begin{array}{c}
                    \layer{k_1}{\ell_1}{\tau_{d_1}}\fatsemi (\layer{k_2}{\ell_2}{\tau_{d_2}}\fatsemi t) \longrightarrow \layer{k_2}{\ell_2+1}{\tau_{d_2-1}}\fatsemi (\layer{k_1+1}{\ell_1}{\tau_{d_1-1}}\fatsemi t) \\
                    \text{if}\;\; d_1\ge1,d_2\ge1 \;\;\text{and}\;\; k_2\le k_1\wedge k_2+d_2-d_1\le k_1<k_2+d_2
                \end{array}
            \end{equation}
        \end{subfigure}

        \vspace{-0.5em}
        \hspace{-0.8em}
        \begin{subfigure}{0.55\textwidth}
            \begin{equation}\label{ru:swap:layerfus}\tag{R19}
                \begin{array}{c}
                    \layer{k_1}{\ell_1}{\tau_{d_1}}\fatsemi\layer{k_2}{\ell_2}{\tau_{d_2}} \longrightarrow \layer{k_2}{\ell_1}{\tau_{d_1+d_2}} \\
                    \text{if}\;\; d_1\ge1,d_2\ge1 \;\;\text{and}\;\; k_1=k_2+d_2
                \end{array}
            \end{equation}
        \end{subfigure}

        \vspace{-0.5em}
        \hspace{-0.8em}
        \begin{subfigure}{0.63\textwidth}
            \begin{equation}\label{ru:swap:layerfusctx}\tag{R20}
                \begin{array}{c}
                    \layer{k_1}{\ell_1}{\tau_{d_1}}\fatsemi (\layer{k_2}{\ell_2}{\tau_{d_2}}\fatsemi t) \longrightarrow \layer{k_2}{\ell_1}{\tau_{d_1+d_2}}\fatsemi t \\
                    \text{if}\;\; d_1\ge1,d_2\ge1 \;\;\text{and}\;\; k_1=k_2+d_2
                \end{array}
            \end{equation}
        \end{subfigure}

        \vspace{-0.5em}
        \hspace{-0.8em}
        \begin{subfigure}{0.63\textwidth}
            \begin{equation}\label{ru:swap:layercom}\tag{R21}
                \begin{array}{c}
                    \layer{k_1}{\ell_1}{\tau_{d_1}}\fatsemi\layer{k_2}{\ell_2}{\tau_{d_2}} \longrightarrow \layer{k_2}{\ell_2}{\tau_{d_2}}\fatsemi \layer{k_1}{\ell_1}{\tau_{d_1}} \\
                    \text{if}\;\; d_1\ge1,d_2\ge1 \;\;\text{and}\;\; k_1>k_2+d_2
                \end{array}
            \end{equation}
        \end{subfigure}

        \vspace{-0.5em}
        \hspace{-0.8em}
        \begin{subfigure}{0.74\textwidth}
            \begin{equation}\label{ru:swap:layercomctx}\tag{R22}
                \begin{array}{c}
                    \layer{k_1}{\ell_1}{\tau_{d_1}}\fatsemi (\layer{k_2}{\ell_2}{\tau_{d_2}}\fatsemi t) \longrightarrow \layer{k_1}{\ell_1}{\tau_{d_1}}\fatsemi (\layer{k_2}{\ell_2}{\tau_{d_2}}\fatsemi t) \\
                    \text{if}\;\; d_1\ge1,d_2\ge1 \;\;\text{and}\;\; k_1>k_2+d_2
                \end{array}
            \end{equation}
        \end{subfigure}
        \vspace{0.2em}
    \end{center}\end{minipage}}
    \label{fig:permrs}
\end{figure}

\section{simp dans Isabelle}

\subsection{en général}

En Isabelle/HOL la tactique \lstinline[language=Isar]|simp| permet de modifier (ou prouver) un but en appliquant des règles de réécriture. En particulier, \lstinline[language=Isar]|simp| peut être invoquée sous plusieurs formes:
\begin{itemize}
    \item \lstinline[language=Isar]|simp| utilise un ensemble de règles de réécriture par défaut. Cet ensemble est un mix de règles judicieusement choisies par les créateurs de bibliothèques et de règles automatiquement générées par des mots clés comme \lstinline[language=Isar]|function|, \lstinline[language=Isar]|fun|, \lstinline[language=Isar]|inductive|, etc.
    \item \lstinline[language=Isar]|(simp add: regle1 regle2)| permet à l'utilisateur de rajouter ses propres règles de réécriture dans l'ensemble utilisé par \lstinline[language=Isar]|simp|
    \item \lstinline[language=Isar]|(simp only: regle1 regle2)| permet à l'utilisateur de ne pas utiliser les règles par défaut, mais seulement celles qu'il choisit
    \item \lstinline[language=Isar]|(simp cong: regle1 regle2)| permet à l'utilisateur de rajouter des règles de congruence
\end{itemize}

Les options \lstinline[language=Isar]|add|, \lstinline[language=Isar]|only| et \lstinline[language=Isar]|cong| peuvent être utilisées dans la même invocation de \lstinline[language=Isar]|simp|, par exemple \lstinline[language=Isar]|(simp only: regle1 cong: regle2)| n'utilise que 2 règles, dont une de congruence.

Une règle de réécriture est un lemme (qui doit être démontré dans Isabelle) sous la forme $P_1 \implies \cdots \implies P_n \implies A = B$. Il représente la règle $A \to B \text{ si } P_1, \ldots, P_n$

En pratique, le simplificateur va essayer de prover les conditions $P_1, \ldots, P_n$ en s'appelant récursivement sur chacune des conditions, et réécrire $A$ en $B$ une fois qu'il a réussi. S'il n'arrive pas à prouver une des conditions (parce qu'il n'a pas assez de règles ou bien parce que la condition est fausse), alors il abandonne l'utilisation de cette règle et en essaye d'autres.

Une règle de congruence est un lemme de la forme $P_1 \implies \cdots \implies P_n \implies f(x) = f(y)$. Elle s'applique avant les règles de réécriture, dès que le simplificateur reconnaît que les deux côtés de l'égalité commencent par $f$. De plus, les règles de congruence se distinguent des règles de réécriture par le fait qu'il ne s'agit pas ici de remplacer $f(x)$ par $f(y)$, mais de prouver $f(x) = f(y)$.

Par exemple, si $f$ est une fonction paire, on peut avoir la règle de congruence $x = -y \implies f(x) = f(y)$, ce qui permet au simplificateur de démontrer $f(-3) = f(3)$

Si le simplificateur arrive à un terme sur lequel il n'y a pas de règles applicables (ou les seules règles potentiellement applicables ont une condition dont la preuve a échoué), alors il s'arrête. Si de plus il s'arrête sur un terme de la forme \verb|t = t|, alors il considère son but comme démontré.


Pour la certification d'équivalence de circuits, il est judicieux d'utiliser \verb|simp|, car notre forme normale est calculée à l'aide de règles de réécriture. Cependant, les règles de réécriture indiquées ci-dessus ne sont pas suffisantes pour le bon fonctionnement de \verb|simp|. En effet, il lui faut aussi:

\begin{itemize}
    \item des règles pour évaluer certaines fonctions ($\dom, \cod, \ldots$)
    \item des règles pour faire de l'arithmétique sur les entiers naturels
    \item des règles pour traiter des inégalités sur les entiers naturels
    \item des règles pour évaluer des opérations booléennes
\end{itemize}

Il est aussi nécessaire d'écrire les règles d'une manière appropriée pour \verb|simp|: une règle comme $f(n+1) = g(n)$ ne s'applique pas sur $f(13)$, mais sur $f(12+1)$. Il convient de prendre à la place une règle comme $n \ge 1 \implies f(n) = g(n-1)$

Il est également important de prendre en compte la difficulté posée par la soustraction sur les entiers naturels. En effet, si $n \ge m$, $m - n$ est définit comme étant $0$ Par exemple, \verb|1 - 2 = 0| peut être démontré en Isabelle. Il faut donc exprimer $n - m$ par \verb|nat (int n - int m)| (on passe par les entiers relatifs)


Un autre obstacle rencontré est le fait que on a besoin de deux règles de congruence:
\begin{equation}
    \nf(u_1) = \nf(v_1) \implies \nf(u_2) = \nf(v_2) \implies \nf(u_1 \otimes u_2) = \nf(v_1 \otimes v_2)
    \label{cong:par}
    \tag{par\_cg}
\end{equation}

\begin{equation}
    \nf(u_1) = \nf(v_1) \implies \nf(u_2) = \nf(v_2) \implies \nf(u_1 \fatsemi u_2) = \nf(v_1 \fatsemi v_2)
    \label{cong:seq}
    \tag{seq\_cg}
\end{equation}

Cependant, Isabelle refusait d'appliquer les deux règles dans le même appel à simp (je pense que c'est parce que le simplificateur ne regarde que la tête du terme, et ne prend qu'une seule règle de congruence par tête, or les deux termes commencent par $\nf$). En effet, un appel à \lstinline[language=Isar]|(simp [...] cong: par_cg seq_cg)| n'applique que la règle \lstinline[language=Isar]|par_cg|, et un appel à \lstinline[language=Isar]|(simp [...] cong: seq_cg par_cg)| n'applique que la règle \lstinline[language=Isar]|seq_cg|


J'ai contourné ce problème en utilisant des \verb|congproc| (ou procédures de congruence). Il est alors possible de donner un pattern plus sophistiqué pour chaque règle.

Par exemple, pour la règle \verb|seq_cg|, j'ai écrit la règle


\begin{lstlisting}[language=Isar]{}
simproc_setup passive weak_congproc cong_seq (\<open>NF (Seq x y)\<close>) =
    \<open>(K (K (K (SOME @{thm seq_cg [cong_format]}))))\<close>
\end{lstlisting}

Dans le code ci-dessus, une \verb|simproc| (procédure de simplification) nommée \verb|cong_seq| est définie. Le mot clé \lstinline[language=Isar]|weak_congproc| sert à déclarer que cette procédure est en réalité une procédure de congruence, et la partie \verb|weak| signale à isabelle que cette procédure ne s'occupe pas de simplifier l'intérieur du terme.

La partie \lstinline[language=Isar]|\<open>NF (Seq x y)\<close>| indique à Isabelle que cette procédure doit être invoquée lorsque le simplificateur essaie de prouver \lstinline[language=Isar]|NF (Seq x1 x2) = NF (Seq y1 y2)|. Cela permet notammment de la distinguer de l'autre règle qui doit être appliquée sur \lstinline[language=Isar]|\<open>NF (Par x y)\<close>|

L'expression \verb|(K (K (K (SOME @{thm seq_cg [cong_format]}))))| est l'équivalent d'une fonction OCaml \verb|fun _ -> fun _ -> fun _ -> Some seq_cq| qui renvoie le théorème correspondant à la règle de congruence voulue. Il est possible de renvoyer un théorème qui dépend des trois arguments, mais ce n'est ici pas nécessaire.

\subsection{avec les PRO}

Pour la simplification des $\pro{\Sigma}$, on utilise l'ensemble de règles de simplification suivant:

\begin{lstlisting}[language=Isar]{}
lemmas PRO_simps =
  R1_NF R2_NF R3_NF R4_NF R5_NF R6_NF R7_NF R8_NF R9_NF R10_NF R11_NF

  admissible.intros              (* rules for admissibility *)
  dom.simps cod.simps def.simps  (* rules for dom, cod & def *)
  trm.disc(5-8)                  (* rules for is_Id *)

  arith_simps                       (* rules for arithmetic *)
  of_nat_0 of_nat_1 of_nat_numeral  (* rules for int *)
  nat_0 nat_numeral                 (* rules for nat *)
  num1                              (* rule to handle 1 *)

  rel_simps(1-71)                (* rules for integer inequalities *)
  simp_thms                      (* rules for boolean calculus *)
\end{lstlisting}

où les règles \verb|R1_NF| jusqu'à \verb|R11_NF| sont les règles de réduction, et les autres règles permettent l'évaluation de fonction et le calcul sur les entiers et les booléens

J'ai également écrit une preuve pour certifier l'équivalence entre deux termes $t_1$ et $t_2$:

\begin{lstlisting}[language=Isar]{}
lemma \<open>NF t1 = NF t2\<close>
proof -
  have pre_eq: \<open>NF (preproc t1) = NF (preproc t2)\<close>
    apply (simp only: preproc.simps)
    including no_limit by (simp only: PRO_simps [[simproc cong_par cong_seq]])?
  have t1_pre: \<open>NF (preproc t1) = NF t1\<close> using pre by blast
  have t2_pre: \<open>NF (preproc t2) = NF t2\<close> using pre by blast
  show ?thesis
    using t1_pre t2_pre pre_eq by argo

qed
\end{lstlisting}

Plusieurs détails sont ici importants:

\begin{itemize}
    \item La fonction \lstinline[language=Isar]|preproc| rajoute quelques circuits vides ($\textup{id}_0$) pour éviter quelques cas problématiques
    \item \lstinline[language=Isar]|including no_limit| permet à la tactique \lstinline[language=Isar]|simp| d'aller à profondeur très grande dans le résultat (pour que \lstinline[language=Isar]|simp| n'abandonne pas sa simplification)
    \item Le point d'interrogation présent à la fin de la ligne 5 permet de gérer correctement le cas où les termes $t_1$ et $t_2$ sont exactement égaux. Dans ce cas, le but est déjà résolu à la ligne 4, et la ligne 5 n'est plus nécessaire. Le problème est donc qu'à la ligne 5, Isabelle
\end{itemize}

\subsection{avec les Perm}

De même, pour les Perm, on utilise cet ensemble de règles:

\begin{lstlisting}[language=Isar]{}
lemmas PRO_Perm_simps =
  R1_to_perm  R2_to_perm  R3_to_perm  R4_to_perm  R5_to_perm
  R6_to_perm  R7_to_perm  R8_to_perm  R9_to_perm  R12_to_perm
  R13_to_perm R14_to_perm R15_to_perm_simp  R16_to_perm_simp
  R17_to_perm_simp  R18_to_perm_simp  R19_to_perm_simp
  R20_to_perm_simp  R21_to_perm_simp  R22_to_perm_simp

  sadmissible.intros             (* rules for admissibility *)
  dom.simps cod.simps def.simps  (* rules for dom, cod & def *)
  trm.disc                       (* rules for is_Id *)

  arith_simps                       (* rules for arithmetic *)
  of_nat_0 of_nat_1 of_nat_numeral  (* rules for int *)
  nat_0 nat_numeral                 (* rules for nat *)
  num1                              (* rule to handle 1 *)

  rel_simps                      (* rules for integer inequalities *)
  simp_thms                      (* rules for boolean calculus *)
\end{lstlisting}

où les règles \verb|R1_to_perm| jusqu'à \verb|R22_to_perm_simp| sont les règles de réduction du système, et les autres règles permettent l'évaluation de fonction et le calcul sur les entiers et les booléens

Et pour certifier l'équivalence entre deux termes $t_1$ et $t_2$:

\begin{lstlisting}[language=Isar]{}
lemma \<open>CF \<lbrakk>t1\<rbrakk> = CF \<lbrakk>t2\<rbrakk>\<close>
proof -
  have pre_eq: \<open>\<lbrakk>preproc t1\<rbrakk> = \<lbrakk>preproc t2\<rbrakk>\<close>
    apply (simp only: preproc.simps)
    including no_limit by (simp only: PRO_Perm_simps [[simproc cong_perm_par cong_perm_seq]])?
  have \<open>sadmissible t1\<close> including no_limit by (simp only: PRO_Perm_simps)
  also have \<open>sadmissible t2\<close> including no_limit by (simp only: PRO_Perm_simps)
  ultimately have \<open>\<lbrakk>t1\<rbrakk> = \<lbrakk>t2\<rbrakk>\<close>
    using pre_eq perm_preproc[of t1] perm_preproc[of t2] by argo
  then show ?thesis by argo

qed
\end{lstlisting}

\section{benchmark}

\subsection{Méthodologie}

Pour tester la rapidité de certification de circuits, des circuits ont été générés. Pour cela un script qui génère des paires de circuits equivalents a été utilisé. Les paires ainsi générées sont ensuite classés par taille et regroupés en lots de 50 en 50 (c'est-à-dire un lot pour les tailles de 0 à 49, un pour les tailles de 50 à 99, etc...).

La taille est définie par le nombre de $\otimes$ et de $\fatsemi$ dans les 2 circuits.

Ensuite, sont générés pour chaque lot de taille un ensemble de 50 exemples en les tirant au hasard (avec remplacement).

Isabelle est ensuite lancé sur les 50 exemples en même temps (ce qui permet à Isabelle de potentiellement paralléliser ses calculs) et le temps final d'éxécution est récupéré.

Le script permettant de générer les circuits est un script aléatoire et ne permet pas de contrôler la taille de sortie des circuits. De plus, pour les très grands circuits, le script peut prendre beaucoup de temps. Pour cette raison, le nombre d'exemples générés dépend de la taille.

Un des problèmes de cette méthode est que la précision de la mesure dépend beaucoup du nombre d'exemples générés: s'il y en a plus de 200, la répartition est proche de l'uniforme, mais lorsqu'il y en a moins que 10, des biais peuvent exister. Il faut donc s'attendre à ce que les tests soient moins précis lorque les circuits sont plus grands (les exemples étant plus longs à générer, on en possède moins)

Un autre problème qui perturbe la mesure est que Isabelle a besoin de s'initialiser et de lire les fichiers qui établissent la théorie permettant de faire la preuve. Ceci prend un temps constant (car ces fichiers ne dépendent pas du circuit testé) qui devrait être assez peu significatif, puisqu'il n'y a qu'une seule étape d'initialisation pour les 50 exemples.

Il est également important de noter qu'Isabelle a pu tourner en parallèle sur 8 threads, ce qui lui permet de traiter plusieurs circuits en même temps.

Tous les exemples ont été lancés sur [TODO insérer nom ordi], effectués sans qu'il y ait d'autres processus en même temps, et tous les tests appartenant à une même famille ont été réalisés durant la même session.

\subsection{PRO bench}

Le nombre d'exemples générés chute très rapidement à partir de la taille 500, du fait du temps de génération.

Le graphe ci-dessous montre le nombre d'exemples générés en fonction de la catégorie. Pour aider à la lisibilité, l'axe des ordonnées est ici logarithmique.

\begin{tikzpicture}
    \begin{semilogyaxis}[
        title={Nombre d'exemples générés par catégorie},
        xlabel={Catégorie},
        ylabel={Nombre d'exemples générés},
        xmin=-20, xmax=620,
        ymin=0.7, ymax=400,
        xtick={0,50,100,150,200,250,300,350,400,450,500,550,600},
        ytick={1,10,100},
        % legend pos=north west,
        ymajorgrids=true,
        grid style=dashed,
        height=7.5cm,
        width=0.9\linewidth
    ]

    \addplot[
        color=blue,
        mark=square,
    ]
    coordinates {
        (0,182)(50,200)(100,221)(150,196)(200,176)(250,130)(300,112)
        (350,105)(400,66)(450,71)(500,23)(550,4)(600,1)
    }
    [pos segment=0,yshift=10pt,font=\footnotesize]
        node [pos=0/12] {182}
        node [pos=1/12] {200}
        node [pos=2/12] {221}
        node [pos=3/12] {196}
        node [pos=4/12] {176}
        node [pos=5/12] {130}
        node [pos=6/12] {112}
        node [pos=7/12] {105}
        node [pos=8/12] {66}
        node [pos=9/12] {71}
        node [pos=10/12] {23}
        node [pos=11/12] {4}
        node [pos=12/12] {1};
        
    \end{semilogyaxis}
\end{tikzpicture}

Le temps de calcul est ici également montré avec l'ordonnée logarithmique. Il est important de noter que la mesure pour la catégorie 600 est probablement imprécise, car on ne dispose que d'un seul exemple (le test s'est donc déroulé avec 50 copies du même exemple, ce qui peut ne pas être représentatif d'un échantillon uniforme).

\begin{tikzpicture}
    \begin{semilogyaxis}[
        title={Temps de vérification par catégorie},
        xlabel={Catégorie},
        ylabel={Temps de vérification [s]},
        xmin=-20, xmax=620,
        ymin=7, ymax=2000,
        xtick={0,50,100,150,200,250,300,350,400,450,500,550,600},
        ytick={10,100,1000},
        % legend pos=north west,
        ymajorgrids=true,
        grid style=dashed,
        height=7.5cm,
        width=0.9\linewidth
    ]

    \addplot[
        color=blue,
        mark=square,
    ]
    coordinates {
        (0, 20.497)(50, 24.821)(100, 30.193)(150, 41.983)(200, 55.789)
        (250, 81.608)(300, 109.642)(350, 120.251)(400, 169.195)
        (450, 231.104)(500, 374.404)(550, 465.433)(600, 1630.060)
    };
        
    \end{semilogyaxis}
\end{tikzpicture}

\subsection{Perm bench}

Ici le nombre d'exemples chute pour les catégories supérieures à 650. Ce nombre est plus grand car il est plus simple de générer des circuits de type Perm que des circuits de type PRO.

Comme le graphe précédent, les ordonnées sont ici logarithmiques.

\begin{tikzpicture}
    \begin{semilogyaxis}[
        title={Nombre d'exemples générés par catégorie},
        xlabel={Catégorie},
        ylabel={Nombre d'exemples générés},
        xmin=-20, xmax=770,
        ymin=0.7, ymax=500,
        xtick={0,50,100,150,200,250,300,350,400,450,500,550,600,650,700,750},
        ytick={1,10,100},
        % legend pos=north west,
        ymajorgrids=true,
        grid style=dashed,
        height=7.5cm,
        width=0.9\linewidth
    ]

    \addplot[
        color=blue,
        mark=square,
    ]
    coordinates {
        (0,250)(50,235)(100,197)(150,205)(200,197)(250,132)(300,123)(350,104)
        (400,83)(450,81)(500,56)(550,35)(600,21)(650,4)(700,2)(750,1)
    }
    [pos segment=0,yshift=10pt,font=\footnotesize]
        node [pos=0/15] {250}
        node [pos=1/15] {235}
        node [pos=2/15] {197}
        node [pos=3/15] {205}
        node [pos=4/15] {197}
        node [pos=5/15] {132}
        node [pos=6/15] {123}
        node [pos=7/15] {104}
        node [pos=8/15] {83}
        node [pos=9/15] {81}
        node [pos=10/15] {56}
        node [pos=11/15] {35}
        node [pos=12/15] {21}
        node [pos=13/15] {4}
        node [pos=14/15] {2}
        node [pos=15/15] {1};
        
    \end{semilogyaxis}
\end{tikzpicture}

Le temps de calcul pour les Perm semble ne pas croîte aussi vite que celui pour les PRO. En conséquence, le graphe ci dessous est donné avec l'axe des ordonnées linéaire et non logarithmique.

Il est important de noter que les points pour les catégories 650, 700 et 750 ne sont probablement pas représentatifs d'un échantillion uniforme, vu leur faible nombre d'exemples.

\begin{tikzpicture}
    \begin{axis}[
        title={Temps de vérification par catégorie},
        xlabel={Catégorie},
        ylabel={Temps de vérification [s]},
        xmin=-20, xmax=770,
        ymin=-50, ymax=1050,
        xtick={0,50,100,150,200,250,300,350,400,450,500,550,600,650,700,750},
        ytick={0,100,200,300,400,500,600,700,800,900,1000},
        % legend pos=north west,
        ymajorgrids=true,
        grid style=dashed,
        height=7.5cm,
        width=0.9\linewidth
    ]

    \addplot[
        color=blue,
        mark=square,
    ]
    coordinates {
        (0, 16.950)(50, 80.245)(100, 65.545)(150, 146.935)
        (200, 241.864)(250, 267.352)(300, 296.208)(350, 361.194)
        (400, 430.655)(450, 466.028)(500, 557.132)(550, 666.792)
        (600, 828.486)(650, 843.943)(700, 805.531)(750, 989.654)
    };
        
    \end{axis}
\end{tikzpicture}

\section{PROP}

Afin d'espérer pouvoir obtenir un résultat similaire sur les $\prop{\Sigma, \varnothing}$, on veut définir une forme normale.

\subsection{Numérotation}

Pour aider à définir une forme normale, on va ajouter des informations à notre circuit.

Cela permet de donner des informations globales sur le circuit de manière locale.

\begin{remark}
    On peut voir le circuit comme un graphe (plus précisément un DAG)
\end{remark}

\begin{theorem}
    $t_1 \simeq t_2 \leqv G(t_1) = G(t_2)$
\end{theorem}

\begin{corollary}
    Toute quantité qui ne dépend que de $G(t)$ est invariante pour toute la classe d'équivalence de $t$
\end{corollary}

\begin{definition}
    On définit la profondeur $\rho(u)$ d'un générateur $u \in \Sigma$ dans un circuit comme étant le nombre maximal de générateur traversés vers la gauche.

    Autrement dit, si les générateurs précédant $v$ sont $u_1, \ldots u_n$, alors
    \[ \rho(v) = \begin{cases}
        0 & \text{si } n=0\\
        1+\max_{1\le i\le n} \rho(u_i) & \text{sinon}
    \end{cases} \]
\end{definition}

\begin{lemma}
    Comme la notion de "précéder" ne dépend que du graphe sous-jacent, la profondeur ainsi définie est indépendante de l'élément de la classe d'équivalence choisie.
\end{lemma}

\begin{definition}
    On note la "vraie profondeur" $p(v)$ d'une porte (c'est-à-dire un générateur ou bien un swap) par l'équation suivante:
    \[
    p(v) = \begin{cases}
        2 \times \rho(v) + 1 & \text{si } v \text{ est un générateur}\\
        2 \times \max_{1\le i\le n} \rho(u_i) & \text{si } v \text{ est un swap et } (u_i)_i \text{ les générateurs le précédant}
    \end{cases}
    \]
\end{definition}

\begin{definition}
    On définit la priorité d'un fil en faisant un parcours en profondeur sur le graphe. On obtient ainsi un ordre topologique sur les fils.
\end{definition}

\begin{lemma}
    Comme cet ordre ne dépend que du graphe sous-jacent, il est équivalent pour tout élément de la classe d'équivalence.
\end{lemma}

\begin{lemma}
    On peut donc munir un circuit des profondeurs et des priorités sans toutefois modifier les classes d'équivalence.
\end{lemma}

\begin{definition}
    On dit d'un circuit (annoté) est en forme normale si et seulement si:
    \begin{itemize}
        \item il est "découpé en couches" selon la profondeur $p$
        \item au sein de chaque couche, les portes sont organisées selon la profondeur de leurs fils de sortie
    \end{itemize}
\end{definition}

\begin{lemma}
    La forme normale existe, et est unique.
\end{lemma}
\begin{proof}
    Existence par correction (DFS) et construction

    Unicité par ordre total des priorités
\end{proof}

\begin{definition}
    On note $\nf$ la procédure qui transforme t en forme normale.
\end{definition}

\begin{property}
    $\nf(\nf(t)) = \nf(t)$
\end{property}

\begin{lemma}
    $t_1 \simeq t_2 \leqv \nf(t_1) = \nf(t_2)$
\end{lemma}
\begin{proof}
    $NF$ ne dépend que du graphe sous-jacent et est unique.
\end{proof}

\begin{remark}
    Comme les $\prop{\Sigma}$ sont des $\pro{\Sigma \cup \Sigma_0}$ (avec des générateurs supplémentaires), on peut tout à fait appliquer la fonction de normalisation des $\pro{\Sigma \cup \Sigma_0}$ sur les $\prop{\Sigma}$, ce qui donne un circuit séparé en couches.

    Celà revient, en outre, à considérer que les swap sont des générateurs quelconques (des boîtes noires, en quelque sorte) et à effectuer le processus de normalisation des $\pro{\Sigma \cup \Sigma_0}$ dessus.
\end{remark}

\begin{definition}
    On pose le système de réécriture suivant sur les circuits numérotés (celui intuitif à deux règles)
\end{definition}

\begin{lemma}
    Tout circuit qui n'est pas en forme normale possède au moins une règle applicable
\end{lemma}

\begin{lemma}
    Le système est terminant
\end{lemma}

\begin{lemma}
    Le système est confluent
\end{lemma}

\begin{remark}
    La règle 1 (échange vertical) peut être remplacée par 3 règles
\end{remark}

\begin{remark}
    La règle 2 (échange à travers une permutation) peut être remplacée par 3 règles (échanges horizontaux; échange à travers un toboggan)
\end{remark}

\begin{definition}
    On pose le système de réécriture (formel) sur les circuits numérotés

    (on met aussi les règles R14 à R22 des Perms)
\end{definition}


\end{document}
